\chapter{Ход работы}
\label{ch:chap2}

\section{СМО М/М/1 без резервирования}

Построим модель простой СМО вида М/М/1 с вероятностью отказа в обслуживании заявки $P_{otk}$. 
Проанализируем влияние потока заявок $\lambda$ на статистические характеристики очереди и сервера.

\begin{figure}[H]
    \centering
    \includegraphics[width=1.0\textwidth]{smo1.png}
    \caption{СМО без резервного сервера}
\end{figure}

\subsection{$\mu > \lambda$}

Параметры модели: $\mu = 6$, $\lambda = 5$, $P_{otk} = 0.5$

\begin{figure}[H]
    \centering
    \includegraphics[width=1.0\textwidth]{base_1.png}
    \caption{Статистические характеристики модели}
\end{figure}

\begin{figure}[H]
    \centering
    \includegraphics[width=1.0\textwidth]{l1.png}
    \caption{Статистические характеристики модели}
\end{figure}

\begin{figure}[H]
    \centering
    \includegraphics[width=1.0\textwidth]{w1.png}
    \caption{Статистические характеристики модели}
\end{figure}

Здесь d - кол-во заявок, вышедших из генератора, ds - кол-во заявок, вышедших из сервера, util - нагрузка на сервер, 
l - средняя длина очереди, w - среднее время ожидания \\

Сравним полученные экспериментальным путем результаты с теоретическими расчетами:

$$L_q = \frac{\lambda^2}{\mu(\mu - \lambda)} = \frac{25}{6} \approx 4.16$$
$$W_q = \frac{\lambda}{\mu(\mu - \lambda)} = \frac{5}{6} \approx 0.83$$

Здесь $L_q$ - средняя длина очереди, $W_q$ - среднее время ожидания. \\

Теоретические значения не сходятся с экспериментальными. Это связано с тем, что формулы не учитывают множество факторов: алгоритм
работы очередей, вместительность очередей. Ко всему прочему в СМО очередей несколько: одна в самой очереди, другая - в сервере. 

\subsection{$\mu < \lambda$}

Параметры модели: $\mu = 6$, $\lambda = 9$, $P_{otk} = 0.5$

\begin{figure}[H]
    \centering
    \includegraphics[width=1.0\textwidth]{base_2.png}
    \caption{Статистические характеристики модели}
\end{figure}

\begin{figure}[H]
    \centering
    \includegraphics[width=1.0\textwidth]{l2.png}
    \caption{Статистические характеристики модели}
\end{figure}

\begin{figure}[H]
    \centering
    \includegraphics[width=1.0\textwidth]{w2.png}
    \caption{Статистические характеристики модели}
\end{figure}

$$L_q = \frac{\lambda^2}{\mu(\mu - \lambda)} = \frac{81}{-54} \approx -1.5$$
$$W_q = \frac{\lambda}{\mu(\mu - \lambda)} = \frac{9}{-54} \approx -0.16$$

Можем заметить, что формулы не работают в случае $\lambda > \mu$ \\

Также можно заметить, что при увеличении $\lambda$ растут параметры $w$, $l$, util. Все это связно с тем, что заявки поступают в СМО
быстрее, чем та может их обработать.

\section{СМО М/М/1 с резервированием}

Резервирование заключается в добавлении в модель еще одного сервера, который бы обрабатывал отброшенные заявки \\

Шаги по добавлению резервирования:

\begin{enumerate}
    \item Добавить новый сервер.
    \item Выход №2 мультиплексора после главного сервера подключить к резервному серверу.
    \item Выход №1 резервного сервера смультиплексировать с выходом №1 главного сервера для вывода заявки из системы
    \item Выход №2 мультиплексора после резервного сервера подключить по обратной связи к очереди
\end{enumerate}

\begin{figure}[H]
    \centering
    \includegraphics[width=1.0\textwidth]{smo2.png}
    \caption{СМО с резервным сервером}
\end{figure}

\subsection{$\mu > \lambda$}

Параметры модели: $\mu = 6$, $\lambda = 5$, $P_{otk} = 0.5$

\begin{figure}[H]
    \centering
    \includegraphics[width=1.0\textwidth]{nn2.png}
    \caption{Статистические характеристики модели}
\end{figure}

Заметим, что параметры модели, которые мы задали, никак не отличаются от парамтеров модели без резервоного сервера. При этом значения
средней длины очереди, среднего времени ожидания и нагрузки сервера значительно уменьшились, что свидетельствует о том, что такая
модель работает быстрее, благодаря резервному серверу, который забирает на себя часть нагрузки, разгружая тем самым очередь и главный
сервер.

\subsection{$\mu < \lambda$}

Параметры модели: $\mu = 6$, $\lambda = 9$, $P_{otk} = 0.5$

\begin{figure}[H]
    \centering
    \includegraphics[width=1.0\textwidth]{nn3.png}
    \caption{Статистические характеристики модели}
\end{figure}

После повышения $\lambda$ нагрузка кратно выросла, но главный сервер нагружен при этом всего лишь на 60-70\%, а в очереди
почти не появляется ожидающих заявок.

\endinput
